% !TeX root = ../papers/nc-p3.tex
\section{特異点を含む整七次点曲線の全係数}
\label{sec:integral-septic-boundary}
\begin{proposition}\label{prop:automatic-spatial-tangent-mark}
四辺の左右が点ファイバーの標準有限Frobenius候補では、共通点グラフの接次元3の全ての点は
標準四移動の固定点であり、それぞれ命題\ref{prop:ore-mark-recognition}のOre表示を与える。
従ってこのような点が一つでもあれば、全次数AS性は三生成部分の有限判定に帰着する。
\end{proposition}
\begin{proof}
任意の接次元3の点からグラフに沿って全位置に点 $p_i$ を取り、初めは周期性を仮定しない。
$H_i=\ker p_i$ とし、定理\ref{thm:maximal-tangent-ore}の証明と同じく個別の評価階数3と結合階数3から、
三つの水平関係 $S_i$ と、残りの文字を右へ移す三つの混合関係を得る。
水平三次交叉の次元を $b$ とすれば全三次元数が $19+b$ 以下なので $b\ge1$。
第一を $p_i$ で縮約して零になる三次交叉は、高さ4の完全対合と一意な前点から
$w(p_{i-1})$ の張る直線である（標準四窓の文字同定を使用）。これが水平三次交叉である。
末端支持は直前の点列から $H_{i-2}$ 全体であり、一方水平性から $H_{i+2}$ に入る。
従って両者は一致する。全位置で同じ議論が通るので四周期であり、Ore認識命題の条件が全て成立する。
\end{proof}

\begin{theorem}\label{thm:integral-septic-boundary}
点ファイバーを持ち点曲線が整な被約七次となる全標準AS窓は、定理\ref{thm:septic-bielliptic-forms}の
八係数のうち、表示した三式が\emph{整な完全交叉曲線}を定めるもの全部である。
全係数で同じ倍率と交換の同型判定、$\lambda\mapsto-\lambda$ による最小周期判定が成り立つ。
この全層は既約で同型類は五次元、全ての特異曲線の窓は滑らかな七次窓から平坦に特殊化する。
\end{theorem}
\begin{proof}[非Gorenstein点の除外]
純性と算術種数5は定理\ref{thm:septic-point-graph}による。どの埋込みにも二次式がない。
滑らかな二次曲面の二次数では種数5にならない。二次錐でも、次数7の整曲線のWeil類は
母線 $A$ と超平面 $H$ により $A+3H$ である。実際解消の負切断を法とした類群は母線が生成し、$H=2A$。
従って反射的な曲線イデアルは $I_{A/Q}(-3)$ と同型で、
$\chi(\mathcal O_C(n))=(n+1)^2-(n-2)(n-3)=7n-5$ となり算術種数6である。
平面七次の種数も5に等しくない。

$h^0(L_i)=4$ である。これ以上なら完全系の五切断による非退化七次の射影四空間模型が得られ、
埋込み四切断を含めるので算術種数5のままであり、Castelnuovo上界3に反する。
双対性により $L_i\hookrightarrow\omega_C$ はスカラーを除き一意で、商は長さ1である。
従って非Gorenstein点があればただ一点 $P$ で、全ての商がそこに支持される。
$P$ の埋込み次元は3である（接次元2以下の被約曲線特異点は平面超曲面または滑らか）。
その極大イデアルを $\mathfrak m$ とすると $\mathfrak m\omega_C\subset L_i$ である。
$\omega_C$ が局所巡回でなく、$\omega_C/L_i$ の長さが1、$L_i$ が自由階数1なので、
$\mathfrak m\omega_C=\mathfrak m L_i$ でなければならない。
実際 $\mathfrak mL_i\subset\mathfrak m\omega_C\subset L_i$ の上の包含が等号なら
$\omega_C/\mathfrak m\omega_C$ が一次元になってしまう。

これで $P$ で消える三文字は全ての位置で $H^0(\mathfrak m_P\omega_C)$ という同じ空間になり、
積を曲線の関数体で比較できる。この水平三文字の間に可換二次式はない（七次を消す二次式がない）。
従ってOre表示の三水平関係は三つの交換子で尽きる。命題\ref{prop:automatic-spatial-tangent-mark}で
混合関係の一次係数により水平空間を順に同定すると、一定の三生成基底の点グラフは平面の射影変換である。
非恒等の場合は定理\ref{thm:plane-ore-unmarked}により二次曲面か平面が四文字の点グラフに入る。
恒等なら一般の持ち上がる水平点で二位置の高さを同時に変えられ、水平射影は曲線を描くので
やはり全グラフの次元が2以上となる。これで非Gorenstein点は除かれる。

\smallskip
\noindent\textit{二つの滑らかな射影中心。}
以下 $\omega_C$ は可逆なので $L_i=\omega_C(-p_i)$、$p_i$ は滑らかな点である
（長さ1のCartier因子は正則点）。滑らかな点だけに支持される相異なる二つの次数3以下の
有効Cartier因子が線型同値になることはない。共通部分を除けば基点自由な次数3以下の鉛筆になる。
その $p_i$ を通るファイバーに必要なら点を加えれば $p_i+D$、$\deg D=2$、
$h^0(p_i+D)\ge2$ となる。残余 $D$ が特異点に乗ってもCartier因子である。
Riemann--Rochにより $h^0(L_i(-D))\ge3$ となり、四切断による閉埋込みが長さ2を分離することに反する。
従って偏極の三階差から得る $p_i+2p_{i+3}\sim2p_{i+1}+p_{i+4}$ は因子の等式である。
定理\ref{thm:septic-point-graph}と同じく、二次式の非存在により異なる二点 $p,q$ が交互に現れ、
$p-q$ の位数はちょうど4である。

$\omega_C$ の五つの切断は閉埋込みを与える。$p$ の外では四切断への射影が既に埋込みであり、
$p$ で消えない第五切断と単純に消える $L_0$ 由来の切断が $p$ の接方向も分離する。
標準二次式は $h^0(\omega_C^2)=12$ により少なくとも三つであり、$p$ の三次元余法空間への微分が
単射なのでちょうど三つである。単射性は $p$ を中心とする射影像に二次式がないことと同値である。
これ以降の六積関係の計算、共通三文字の三交換子と順序付き二点からの同型復元は
定理\ref{thm:septic-point-graph}の証明が適用できる。

\smallskip
\noindent\textit{係数の必要性と逆構成。}
定理\ref{thm:septic-bielliptic-forms}の三標準式は整な完全交叉である。
共通三文字 $J$ が非零の開部分はグラフから復元でき、$J=0$ の直線には二つの指定点しかないからである。
二点が滑らかで接線が曲線に含まれないことから、同定理の三次重なり計算で
$xy=D,\ xr-ry=F,\ xs+sy=G$ を得る。ここで $r,s,h$ は共通三文字である。
$D,F,G$ の $\langle r^2,s^2\rangle$ を法とした方向は一本である。
共通代表の $h^2$ 係数が零だとすると $h$ 軸の点で六関係を左評価した像は
$\langle r,s\rangle$ に入る。実際三交換子の評価もここに入り、残る三式には $hx,hy,h^2$ がない。
従って右点ファイバーが少なくとも直線となり矛盾する。
平方完成と二つの補文字の変更により滑らかな場合と同じ八係数表示になる。

逆にこの表示が整な完全交叉を定めれば、両指定点 $r=s=h=0$, $v=\pm u$ は滑らかである
（三式のそこでの微分は $dr,ds,dv\mp du$ の非零倍）。長さ3のスキームを含む直線は三式で
恒等的に消えて曲線成分となるので存在しない。従って各指定点からの完全射影は閉埋込みである
（長さ2の分離が失敗すれば、指定点と合わせた長さ3が直線に乗る）。
$J\ne0$ での点対の復元は座標比の正則な逆を持つため、特異点でもスキーム同型である。
一方だけで $J=0$ の箇所は滑らかな指定点での厳密変換とその三つの独立な微分だけを用いた
前掲定理の計算で同型になる。両側で $J=0$ は $xy$ の式で不可能である。

全次数の計算にも新たな例外はない。$F,G,D$ の $H=h^2+\lambda rs$ 係数の少なくとも一本は非零である
（全て零なら三式は $h$ と独立で、整な非退化曲線にならない）。
従って中央三二次式上の単項式基底と三次重なりの計算が成立する。
双対の可換三二次式に共通射影零点がない理由は整曲線上で二つの非零切断の積が零にならないことであり、
これも特異点を許す。定理\ref{thm:septic-bielliptic-forms}の有限Frobenius計算と
命題\ref{prop:commuting-three-point-graph}が全次数の正則性を与える。
二つの射影中心と二つの固有文字の復元も変わらず、従って同型と周期の倍率・交換判定が全域で成立する。
射影完全交叉の平坦な開集合で幾何学的整性は開であり、滑らかな係数はその中で稠密である。
任意の許容係数を一般の直線で滑らかな係数につなぐと、各次数のHilbert関数の一致から
代数族もその点の離散付値環上で平坦となる。同型類の次元は滑らかな開集合の5である。
\end{proof}

\paragraph{特異な許容係数の例。}
$\lambda=0$, $F=S+H$, $G=R+H$, $4D=3R-3S+8H$ とする。
$r=1$, $s=x\ne0$, $z=h^2$ と置けば $u=x^2+z$, $v=(1+z)/x$、残る式は
\[
        (1-x^2)z^2+(2-2x^4-8x^2)z+(1-x^2)^3=0
\]
である。判別式は $16x^2(1+x^2)(3x^2-1)$ で平方でない。
$x=1,z=0$ では $z$ の係数が非零で解の位数は3、従ってこの二次関数体で $z$ も平方でない。
$h^2=z$ の全曲線はこの開集合で一般点が被約かつ既約であり、原始多項式なので縦の曲線成分もない。
$r=0$ または $s=0$ に三式の解は有限個しかなく、完全交叉の純性により曲線は整である。
$x=1,h=0$ は尖点であり、この窓も全次数で正則となる。
